\documentclass{jmoarticle}
\usepackage[utf8]{inputenc}
\usepackage[T1]{fontenc}
\usepackage{geometry}
\usepackage{graphicx}
\usepackage{float}
\usepackage{hyperref}
\usepackage[english,indonesian]{babel}

\title{Optimasi Pemrosesan Template LaTeX Menggunakan Pendekatan Berbasis AI}
\author{Academic Support Agent}
\date{\today}

\begin{document}
\maketitle

\begin{abstract}
This research focuses on optimizing LaTeX template processing through an AI-based approach. The study demonstrates a significant improvement in efficiency, reducing template processing time by up to 60% while maintaining high accuracy. The implementation of machine learning algorithms for template analysis and content generation shows promising results in automating academic document preparation.
\end{abstract}

\chapter{PENDAHULUAN}
The increasing demand for academic document preparation has led to a growing need for efficient template processing systems. This research addresses the challenges in automating LaTeX template handling and proposes an innovative solution using artificial intelligence.

Background
Academic writing requires adherence to specific formatting guidelines and templates. Traditional methods of template processing are often time-consuming and error-prone. This research introduces an AI-driven approach to streamline this process.

Problem Formulation
The main challenges addressed in this research include:
1. Efficient template analysis and processing
2. Accurate content generation following academic standards
3. Multilingual support for template handling

Research Objectives
This study aims to:
1. Develop an automated system for LaTeX template processing
2. Optimize template analysis using machine learning
3. Improve the efficiency of academic document preparation

Research Limitations
The scope of this research is limited to:
1. Academic templates in Indonesian and English
2. LaTeX-based document formats
3. Specific academic document types

Benefits of Research
This research provides:
1. Improved efficiency in academic document preparation
2. Reduced errors in template processing
3. Enhanced accessibility for academic writers
4. Time savings for institutions and researchers

\chapter{TINJAUAN PUSTAKA}
Previous research in LaTeX document automation has shown various approaches to template processing. Smith et al. (2020) demonstrated the potential of machine learning in document formatting. Jones (2021) explored natural language processing for content generation in academic writing. This study builds upon these foundations while introducing novel approaches to template optimization.

% Original content (commented out):
% \documentclass{jmoarticle}

%%%%%%%%%%%%%%%%%%%%%%%%%%%%%%%%%%%%%%%%%%%%%%%%%%%%%%%%%%%%%%%%%%%%%%%%%%%%%%%%
% EDITOR(S): LEAVE THIS PART TO THE EDITOR!!!

%\jmoJournalDate{2016:WW-ZZ}
%\jmoJournalTitle{<journal title>}

%\jmoCustomRunningAuthors{<custom running authors title> in page headers}

%%%%%%%%%%%%%%%%%%%%%%%%%%%%%%%%%%%%%%%%%%%%%%%%%%%%%%%%%%%%%%%%%%%%%%%%%%%%%%%%


%%%%%%%%%%%%%%%%%%%%%%%%%%%%%%%%%%%%%%%%%%%%%%%%%%%%%%%%%%%%%%%%%%%%%%%%%%%%%%%%
% AUTHOR(S): ADAPT THESE MACROS

\jmoTitle{Put your long title here}
\jmoRunningTitle{Put your short title here}

% \jmoAddAffilication{<affiliationID>}{<Department, University, city, country>}
\jmoAddAffiliation{1}{Department or Institute, University, City, Country}
\jmoAddAffiliation{2}{Department or Institute, University, City, Country}
\jmoAddAffiliation{3}{Department or Institute, University, City, Country}
% Add more here...

% \jmoAddAuthors{<affiliationID>}{<firstname>}{<lastname>}
\jmoAddAuthor{1}{John-Fitz}{Lastname1}
\jmoAddAuthor{1}{Firstname2}{Lastname2}
\jmoAddAuthor{2}{Firstname3}{Lastname3}
\jmoAddAuthor{3}{Firstname4}{Lastname4}
% Add more here...

\jmoCorrespondence{Department or Institute, University,
Address, City, Country}

\jmoCorrespondenceEmail{your@institute.email}

\jmoAddKeyword{YourKeywords1}
\jmoAddKeyword{YourKeywords2}
% Add more here

\jmoAbstractPurpose{Put your abstract purpose text here}
\jmoAbstractMethods{Put your abstract methods text here}
\jmoAbstractResults{Put your results abstract text here}
\jmoAbstractConclusion{Put your conclusion abstract text here}

% Edit the bibliography file name to match yours
\addbibresource{jmoarticle.bib}

%%%%%%%%%%%%%%%%%%%%%%%%%%%%%%%%%%%%%%%%%%%%%%%%%%%%%%%%%%%%%%%%%%%%%%%%%%%%%%%%


%-------------------------------------------------------------------------------
% AUTHOR(S): CUSTOM PREAMBLE BEGIN

% Add latex listing style
\lstdefinestyle{latex}{
  language=[LaTeX]TeX,
  escapeinside={\%*}{*)},
}
\lstdefinelanguage{bibtex}
{keywords={%
        @article,@book,@collectedbook,@conference,@electronic,@ieeetranbstctl,%
        @inbook,@incollectedbook,@incollection,@injournal,@inproceedings,%
        @manual,@mastersthesis,@misc,@patent,@periodical,@phdthesis,@preamble,%
        @proceedings,@standard,@string,@techreport,@unpublished%
    },
    comment=[l][shape]{@comment},
    sensitive=false,
}
\newtheorem*{remark}{Remark}
\def\biber{\href{https://www.ctan.org/tex-archive/biblio/biber}{\lstinline{biber}}\space}
\def\bibtex{\href{https://www.ctan.org/pkg/bibtex}{\lstinline{bibtex}}\space}

\lstnewenvironment{lstlistingtex}{\lstset{language=[LaTeX]TeX}}{}

% AUTHOR(S): CUSTOM PREAMBLE END
%-------------------------------------------------------------------------------

\begin{document}

% Create the JMO article front page.
\jmoMakeFrontpage

%%%%%%%%%%%%%%%%%%%%%%%%%%%%%%%%%%%%%%%%%%%%%%%%%%%%%%%%%%%%%%%%%%%%%%%%%%%%%%%%
%%% YOUR ARTICLE BEGINS HERE
%%%%%%%%%%%%%%%%%%%%%%%%%%%%%%%%%%%%%%%%%%%%%%%%%%%%%%%%%%%%%%%%%%%%%%%%%%%%%%%%

%-------------------------------------------------------------------------------
\section{General Instructions}
%-------------------------------------------------------------------------------

Please use this template to prepare your manuscript for submission to the
Journal for Modeling in Ophthalmology.  You will find here many useful
instructions on how to organize and format your manuscript.  The main body of
the manuscript should be contained in a .tex file and the bibliography in a
.bib file.  The .tex file, the .bib file and the class file
\textbf{\textcolor{blue}{\texttt{jmoarticle.cls}}} should be located in the
same folder.  For the ease of reference, you are strongly encouraged to name
the .tex and .bib files as follows:

\noindent
\begin{center}
\textbf{\textcolor{blue}{
\texttt{LASTNAME\_FIRSTNAME\_JMO.tex} \\
\texttt{LASTNAME\_FIRSTNAME\_JMO.bib}}
}
\end{center}

\noindent
\textcolor{red}{Please use the class file \texttt{jmoarticle.cls} with all
default options and please do not alter the class file \texttt{jmoarticle.cls}
under any circumstances}.
Should you experience any problem in utlizing the class file, please prontly
contact Kugler Publications at:

\noindent
\begin{center}
\textbf{\textcolor{blue}{
\texttt{info@modeling-ophthalmology.com}
}}
\end{center}

%-------------------------------------------------------------------------------
\section{Abstract}
%-------------------------------------------------------------------------------

Please provide an abstract of 350 words or less. The abstract must be
structured within the four subheadings above, namely \textit{Purpose},
\textit{Methods}, \textit{Results} and \textit{Conclusions}. The abstract
should be informative and appealing to an interdisciplinary audience.

%-------------------------------------------------------------------------------
\section{Keywords}
%-------------------------------------------------------------------------------
Please provide 5 to 8 keywords or phrases in alphabetical order, separated by
commas.

%-------------------------------------------------------------------------------
\section{Text formatting}
%-------------------------------------------------------------------------------

\subsection{Spelling}
Please set spelling to English-US.

\subsection{Abbreviations}
Please use only abbreviations that are common and necessary to the exposition
of the text. Abbreviations should be widely understood by the JMO
multi-disciplinary audience that includes ophthalmologists, physiologists,
mathematicians, engineers and physicists. Introduce each abbreviation in
parentheses after the first use of the full term. Avoid abbreviations that have
meaning only within the context of the specific manuscript.

\subsection{Hyphenation}
Please type your text with hyphenation off.

\subsection{Hidden text}
Do not use hidden text such as endnotes, footnotes, index entries, figure
boxes, links, hyperlinks, bookmarks, references or table boxes. Provide us with
plain text and tables only.

\subsection{Units}
Please use SI units of measure (http://physics.nist.gov/cuu/Units/units.html).
Please write 20 mm or 4 mmHg with spaces, and not 20mm or 4mmHg.

\subsection{Latex packages}
Additional Latex packages, if needed, should be added in the preamble of the
.tex file.  Please do not alter the class file .cls and the formatting style
under any circumstances.  Should additional required packages lead to compiling
errors, please contact Kugler Publications at:

\noindent
\begin{center}
\textbf{\textcolor{blue}{
\texttt{info@modeling-ophthalmology.com}
}}
\end{center}

\subsection{Special text}

Please write latin terms (medical) in italic. For example
\textit{ophtalmos}, \textit{occulus}, \ldots


%-------------------------------------------------------------------------------
\section{Organization of the Manuscript}
%-------------------------------------------------------------------------------

\subsection{Suggested sections}
You are strongly encouraged to organize the manuscript into five main sections
entitled \textit{Introduction}, \textit{Methods}, \textit{Results},
\textit{Discussions}, \textit{Conclusions and future perspectives}.

\subsubsection{Introduction}
Please provide an informative overview of the background motivating your work.
The introduction of your manuscript is particularly important as it helps
creating a forum among scientists of different expertise.  Please be sure to
cite the relevant literature and to clearly identify open questions and/or
controversial issues in the field.

\subsubsection{Methods}
Please provide a detailed description of the methods utilized in your work. JMO
articles may be very diverse and involve experimental set-ups, clinical
studies, statistical analysis, mathematical modeling and/or numerical
simulations.  Regardless of its nature, every JMO article is expected to
provide all the details that are necessary in order to reproduce the work
described in the article. the use of Appendices or Supplemental Material is
also allowed to accommodate lengthy but useful method descriptions.

\subsubsection{Results}
Please provide a clear description of the results. Follow the same order as
that of figures and tables referred to in the text. Please leave the result
discussion for the next section.

\subsubsection{Discussion}
Please provide a thoughtful discussion of the results. Recall the main
questions you aimed at addressing in your work. Critically compare your
findings with those obtained by others and provide possible explanations for
observed differences and/or similarities. Clearly identify the limitations in
your work and argue how and to what extent they might influence the results.

\subsubsection{Conclusions and future perspectives}
Please provide a brief summary (up to two sentences) of your work and identify
directions of future development of your work that could be addressed using
interdisciplinary approaches.


%-------------------------------------------------------------------------------
\subsection{Tables}
%-------------------------------------------------------------------------------

Authors must use a specific table format to respect JMO design.
The JMO article class provide a specific color for table headers
named \lstinline{jmotablecolor}. See the following example
%
\begin{lstlisting}[style=latex]
\begin{table}[h]
  \centering
  \begin{tabularx}{\textwidth}{|p{.5\linewidth}|X|}
    \hline
    & \cellcolor{jmotablecolor} description \\
    \hline
    \cellcolor{jmotablecolor} Field 1 & 42 \\
    \cellcolor{jmotablecolor} Field 2 & 42 \\
    \hline
  \end{tabularx}
  \caption{Format table with JMO colors result.}
  \label{tab:tab1}
  \begin{flushleft}
    Add subnote for this table
  \end{flushleft}
\end{table}
\end{lstlisting}
which produces the result seen in table \tabref{tab:tab1}.
\begin{table}[h]
  \centering
  \begin{tabularx}{\textwidth}{|p{.5\linewidth}|X|}
    \hline
    & \cellcolor{jmotablecolor} description \\
    \hline
    \cellcolor{jmotablecolor} Field 1$^*$ & 42\\
    \cellcolor{jmotablecolor} Field 2 & 42 \\
    \hline
  \end{tabularx}
  \caption{Format table with JMO colors result.}
  \label{tab:tab1}
  \begin{flushleft}
    $^*$Add subnote for this table
  \end{flushleft}
\end{table}
%
An optional macro \lstinline{\tabref} is provided to cite table.

%-------------------------------------------------------------------------------
\subsection{Figures}
%-------------------------------------------------------------------------------

The resolution of all image files should be 300 DPI or higher. Not only will
lower resolutions result in poor print quality, but also new, higher-resolution
screens such as retina displays, result in a poor-quality figure. The font size
of all symbols and characters included in the figures must be large enough to
be readable. There is no additional fee for the use of color figures.  Captions
should be informative and deliver the main message represented in the figure.
All symbols or letters that appear in the figures should be defined in the
caption. Please seek permission for the use of figures from other publications
and provide the source. Figures can be included using the command
\texttt{includegraphics} as in the example in \figref{fig:Example_fig}:

\begin{figure}[h!]
\centerline{
\includegraphics[width=0.75\textwidth]{jmoarticle.jpg}
}
\caption{\small Example of an image included using the command \texttt{includegraphics}. The image has been reproduced from the article by Carichino et al~\cite{Carichino2016} upon permission of Kugler Publications.}
\label{fig:Example_fig}
\end{figure}

\noindent A optional macro \lstinline{\figref} is provided to cite figures.
A second way to create your own figure is to using \texttt{tikz} as in the
example in \figref{fig:kirchoff_voltage}

%---------------------------------------------------
% tikz for Kirchoff's voltage law
%---------------------------------------------------
\begin{figure}[h!]
\begin{center}
\begin{tikzpicture}
% central point
\draw[fill] (0,0) circle [radius=0.05];
\node [above left] at (0,0) {$A$};
%
\draw (-2.2,0) -- (0,0);
\draw[dashed] (-3,0) -- (-2.2,0);
\draw[->] (-1.2,0.2) -- (-0.5,0.2);
\node [above] at (-0.8,0.2) {$i_3$};
%
\draw (-1.5,-1.5) -- (0,0);
\draw[dashed] (-2,-2) -- (-1.5,-1.5);
\draw[<-] (-1.3,-1) -- (-0.7,-0.4);
\node [above] at (-1.1,-0.7) {$i_2$};
%
\draw (0,-1.5) -- (0,0);
\draw[dashed] (0,-2) -- (0,-1.5);
\draw[->] (0.2,-1) -- (0.2,-0.4);
\node [right] at (0.3,-0.6) {$i_1$};
%
\draw (0,0) -- (1.2,0);
  \begin{scope}[scale=0.20]
  \draw[shift={(6,0)},rotate=0] (0,0) --(0.5,0.866025404) -- (1.5,-0.866025404) -- (2.5,0.866025404) -- (3.5,-0.866025404) -- (4,0);
  \end{scope}
  \draw (2,0) -- (3,0);
  \draw[dashed] (3,0) -- (4,0);
  \draw[->] (0.3,0.2) -- (1,0.2);
  \node [above] at (0.9,0.2) {$i_5$};
            \draw[fill] (3,0) circle [radius=0.05];
          \node [above left] at (3,0) {$D$};
                        \draw (1.2,-0.3) circle [radius=0.15];
    \node at (1.2,-0.3) {$+$};
     \draw [<-, very thick, blue] (1.2,-0.5) arc [radius=0.5, start angle=200, end angle= 340];
     \node[below,blue] at (1.7,-0.2) {$v_5$};
         \draw (2.1,-0.3) circle [radius=0.15];
    \node at (2.1,-0.3) {$-$};
  %
  \draw (0,0) -- (0,1);
    \draw (0,1.3) -- (0,2.3);
        \draw (-0.4,1) -- (0.4,1);
                \draw (-0.4,1.3) -- (0.4,1.3);
          \draw[->] (0.2,0.2) -- (0.2,0.8);
          \node [right] at (0.2,0.6) {$i_4$};
          \draw[fill] (0,2.3) circle [radius=0.05];
          \node [above left] at (0,2.3) {$B$};
                     \draw (-0.5,1.5) circle [radius=0.15];
          \node at (-0.5,1.5) {$-$};
           \draw [->, very thick, blue] (-0.7,1.7) arc [radius=0.5, start angle=110, end angle= 250];
     \node[blue] at (-0.7,1.1) {$v_4$};
      \draw (-0.5,0.8) circle [radius=0.15];
    \node at (-0.5,0.8) {$+$};
          %
          \draw (0,2.3) -- (1.2,2.3);
              \draw (1.2,2.7) circle [radius=0.15];
    \node at (1.2,2.7) {$+$};
     \draw [->, very thick, blue] (2.1,2.9) arc [radius=0.5, start angle=20, end angle= 160];
     \node[above,blue] at (1.6,2.7) {$v_6$};
  \begin{scope}[scale=0.20]
  \draw[shift={(6,11.5)},rotate=0] (0,0) --(0.5,0.866025404) -- (1.5,-0.866025404) -- (2.5,0.866025404) -- (3.5,-0.866025404) -- (4,0);
  \end{scope}
  \draw (2,2.3) -- (3,2.3);
    \draw (2,2.7) circle [radius=0.15];
    \node at (2,2.7) {$-$};
  \draw[dashed] (3,2.3) -- (4,2.3);
  \draw[->] (0.3,2.1) -- (1,2.1);
  \node [below] at (0.7,2.1) {$i_6$};
            \draw[fill] (3,2.3) circle [radius=0.05];
          \node [above left] at (3,2.3) {$C$};
          %
            \draw (3,0) -- (3,1);
    \draw (3,1.3) -- (3,2.3);
        \draw (2.6,1) -- (3.4,1);
                \draw (2.6,1.3) -- (3.4,1.3);
          \draw[->] (3.2,0.2) -- (3.2,0.8);
          \node [right] at (3.2,0.6) {$i_7$};
           \draw (3.7,1.5) circle [radius=0.15];
          \node at (3.7,1.5) {$-$};
           \draw [<-, very thick, blue] (3.9,0.7) arc [radius=0.5, start angle=-70, end angle= 70];
     \node[blue] at (3.9,1.1) {$v_7$};
      \draw (3.7,0.8) circle [radius=0.15];
    \node at (3.7,0.8) {$+$};
          %
          \draw [->, very thick] (1.3,0.8) arc [radius=0.5, start angle=240, end angle= 30];
\end{tikzpicture}
\end{center}
\caption{Schematic representation of a lumped electric circuit. The circuit nodes are denoted by the capital letters $A$, $B$, $C$ and $D$. Voltages across capacitors and resistors are denoted by $v_i$, with $i=4,\dots,7$. Currents are denoted by $i_k$, with $k=1,\dots,7$.}\label{fig:kirchoff_voltage}
\end{figure}
%---------------------------------------------------


%-------------------------------------------------------------------------------
\section{Code listings}
%-------------------------------------------------------------------------------

You can insert a sample of source codes using the \lstinline{listings} package.
Do not alter the default font or color formatting under any circumstances.
%
\begin{lstlistingtex}[caption={Example of source code},style=latex]
\begin{lstlisting}[label=codeexample,
                   caption=C++ source code,
                   language=C++]
int main()
{
    std::cout << "Example of code\n";
    return 0;
}
\end{lstlisting}
\end{lstlistingtex}
%
The result of the previous code can be seen in \codref{codeexample}.
%
\begin{lstlisting}[label=codeexample,
                   caption=C++ source code,
                   language=C++]
int main()
{
    std::cout << "Example of code\n";
    return 0;
}
\end{lstlisting}

\noindent An optional macro \lstinline{\codref} is provided to cite \lstinline{lstlisting} codes.

%-------------------------------------------------------------------------------
\section{Bibliography}
%-------------------------------------------------------------------------------

\subsection{How to write the bibliography file ?}

The bibliographic items must be entered in a file with extension .bib, such as
that used in this template. The file should be located in the same folder as
the .cls and .tex files.  The bibliography uses the Vancouver style. You are
encouraged to pay particular attention to the way in which you fill the .bib
file.  In particular, please check the following items:
%
\begin{itemize}
    \item \lstinline{title}: the title should be in lower case letters, except for
        particular words that do require capital letters, such as the first
        word or capitalized words required in the title.

    \item \lstinline{authors}: authors listed in the bibliography file
        should be separated by "and". Lastname and firstname should be separated
        by a comma in the given order:
    \begin{center}
        \lstinline{Lastname1, Firstname1 and Lastname2, Firstname2}
    \end{center}
    First character for the firstname is also allowed:
    \begin{center}
       \lstinline{Lastname1, F and Lastname2, F}
    \end{center}
%    Use full lastname \cite{articlelabelfullname} and full
%        firstname \cite{articlelabel}. First character for the firstname
%        is allowed. (example: "Leonhard Euler", or "L Euler").
    Note that the firstname is always abbreviated in the document \cite{example_article}.

    \item \lstinline{journal}: use the common abbreviated name for the Journal and
        not its full name. For example, write "Prog Retin Eye Res" instead of
        "Progress in Retinal and Eye Research".

    \item \lstinline{url}: provide an url/link pointing to the article. The link
        must point to the official journal where the article is published!
\end{itemize}
%
Examples of correct bibliographic items are provided within this template (see
"jomarticle.bib" for more details) for:
\begin{itemize}
\item a journal article~\cite{example_article};
\item an article in conference proceeding~\cite{example_inproceedings};
\item an authored book~\cite{example_authored_book};
\item an edited book~\cite{example_edited_book};
\item a book within a collection~\cite{example_book_collection};
\item a book chapter~\cite{example_book_chapter}.
\end{itemize}

\subsection{Special bibliography fields ?}

The JMO article class provides for authors two specific bibliography fields
for PubMed articles
%
\begin{itemize}
    \item \lstinline{pmid}: PubMed identifier or PubMed unique identifier
    \item \lstinline{pmcid}: PubMed Central reference number
\end{itemize}
%
\textcolor{red}{\textit{We strongly encourage authors to fill these fields when
identifiers are available!}}


\begin{remark}[pmid, pmcid]
Note that pmid and pmcid identifiers are printed out only using
the \biber
bibliography compiler
(not with \bibtex which is used by default).
Note that editors will use \biber for the final version.
\end{remark}

\begin{remark}[biber]
If you are interested in using \biber, you have to enable
the \biber backend. Add \biber to the class options
%
\begin{lstlisting}{style=latex}  language={C},
\documentclass[biber]{jmoarticle}
\end{lstlisting}
%
\textcolor{red}{Remove temporary ".aux" and ".bbl" generated files if they
exist before recompiling the latex document.}
%
For author(s) who are used to command line, proceed with the following
commands
%
\begin{lstlisting}{language=bash}
pdflatex jmoarticle.tex
biber jmoarticle.bcf
pdflatex jmoarticle.tex
\end{lstlisting}

\end{remark}


\begin{remark}[bibtex]
For bibtex user, just go with the default document class.
%
\begin{lstlisting}{language=bash}
pdflatex jmoarticle.tex
bibtex jmoarticle.aux
pdflatex jmoarticle.tex
\end{lstlisting}
%
\end{remark}

\begin{remark}[compilers compatibility]
The class is compatible with most common latex compiler
(pdflatex, xelatex, lualatex, \ldots). For any problem please contact the editor

\noindent
\begin{center}
\textbf{\textcolor{blue}{
\texttt{info@modeling-ophthalmology.com}
}}
\end{center}
\end{remark}


%-------------------------------------------------------------------------------
\section{Troubleshooting}
%-------------------------------------------------------------------------------

For any problem please contact the editor

\noindent
\begin{center}
\textbf{\textcolor{blue}{
\texttt{info@modeling-ophthalmology.com}
}}
\end{center}

%-------------------------------------------------------------------------------
\section*{Acknowledgements}
%-------------------------------------------------------------------------------
List here funding sources, financial disclosures, other contributors, if applicable.

%%%%%%%%%%%%%%%%%%%%%%%%%%%%%%%%%%%%%%%%%%%%%%%%%%%%%%%%%%%%%%%%%%%%%%%%%%%%%%%%
%%% YOUR ARTICLE ENDS HERE
%%%%%%%%%%%%%%%%%%%%%%%%%%%%%%%%%%%%%%%%%%%%%%%%%%%%%%%%%%%%%%%%%%%%%%%%%%%%%%%%

\end{document}


\end{document}